\documentclass[12pt]{article}
\usepackage[utf8]{inputenc}
\usepackage[french]{babel}
\usepackage[T1]{fontenc}
\usepackage{geometry}
\usepackage{xcolor}
\usepackage{enumitem}

\geometry{margin=2cm}
\definecolor{mainblue}{RGB}{0,102,204}
\definecolor{importantred}{RGB}{153,0,0}

\begin{document}

\begin{center}
    \Huge \textbf{\textcolor{mainblue}{Oral : Les Temps Modernes}} \\
    \large \textit{Focus : Arts, Musique et Pouvoir (1492 - 1789)}
\end{center}

\vspace{0.8cm}

\section*{\textcolor{mainblue}{I. La Renaissance : L'Art comme Science}}
\begin{itemize}[label=\textcolor{mainblue}{$\bullet$}]
    \item \textbf{\textcolor{importantred}{Léonard de Vinci :}} 
        \begin{itemize}[label=--]
            \item \textbf{L'Homme de Vitruve :} Symbole de l'humanisme. L'homme est au centre de l'univers, lié aux mathématiques.
            \item \textbf{Sfumato :} Technique vaporeuse (La Joconde) pour un réalisme psychologique.
            \item \textbf{Anatomie :} Disssections pour comprendre la mécanique du corps.
        \end{itemize}
    \item \textbf{\textcolor{importantred}{Musique & Imprimerie :}} 
        \begin{itemize}[label=--]
            \item Diffusion massive des partitions (1450).
            \item Développement de la \textbf{polyphonie} (mélange de plusieurs voix indépendantes).
        \end{itemize}
\end{itemize}

\vspace{0.5cm}

\section*{\textcolor{mainblue}{II. L'Âge Baroque : Le Spectacle du Pouvoir}}
\begin{itemize}[label=\textcolor{mainblue}{$\bullet$}]
    \item \textbf{\textcolor{importantred}{Louis XIV & La Danse :}} 
        \begin{itemize}[label=--]
            \item Le Roi est danseur étoile. La danse est un \textbf{outil politique} pour soumettre la noblesse.
            \item Fondation de l'Académie Royale de Danse (1661) : naissance des codes classiques.
        \end{itemize}
    \item \textbf{\textcolor{importantred}{Les Maîtres du Baroque :}} 
        \begin{itemize}[label=--]
            \item \textbf{Lully :} Créateur de l'Opéra français. Musique imposante pour la gloire du Roi.
            \item \textbf{Vivaldi :} Le \textit{Concerto} moderne. Musique descriptive (Les Quatre Saisons).
            \item \textbf{Bach :} Maître absolu de la \textbf{Fugue}. Il fixe les règles de l'harmonie moderne.
        \end{itemize}
\end{itemize}

\vspace{0.5cm}

\section*{\textcolor{mainblue}{III. Les Lumières : Vers la Modernité}}
\begin{itemize}[label=\textcolor{mainblue}{$\bullet$}]
    \item \textbf{\textcolor{importantred}{Mozart & Le Classicisme :}} 
        \begin{itemize}[label=--]
            \item Recherche de clarté, d'équilibre et de simplicité (fin du XVIIIe).
            \item \textbf{Émancipation :} L'artiste devient indépendant (plus seulement au service d'un prince).
        \end{itemize}
    \item \textbf{\textcolor{importantred}{Évolution Technique :}} 
        \begin{itemize}[label=--]
            \item Transition majeure : du \textbf{Clavecin} (pincé) au \textbf{Piano-forte} (frappé).
            \item Apparition des premiers \textbf{concerts publics} payants.
        \end{itemize}
\end{itemize}

\vspace{0.8cm}

\begin{center}
    \textit{\small Rappel : Garder un ton dynamique et faire le lien avec les visuels des slides Gamma.}
\end{center}

\end{document}